% Generated draft: requires mathematical and source review.
\documentclass[11pt]{article}
\usepackage[margin=25mm]{geometry}
\usepackage{fontspec}
\setmainfont{Latin Modern Roman}
\usepackage{amsmath,amssymb,mathtools,unicode-math}
\setmathfont{Latin Modern Math}
\usepackage{enumitem,needspace,xurl,hyperref}
\setlength{\parindent}{0pt}
\setlength{\parskip}{5pt}
\setlength{\emergencystretch}{3em}
\newcommand{\N}{\mathbb{N}}
\newcommand{\Z}{\mathbb{Z}}
\newcommand{\Q}{\mathbb{Q}}
\newcommand{\R}{\mathbb{R}}
\newcommand{\C}{\mathbb{C}}
\newcommand{\F}{\mathbb{F}}
\newcommand{\A}{\mathbb{A}}
\newcommand{\PP}{\mathbb{P}}
\newcommand{\E}{\mathbb{E}}
\newcommand{\eps}{\varepsilon}
\newcommand{\dd}{\,\mathrm{d}}
\begin{document}
\section*{UnsolvedMath ID: 30002308}
\textbf{Source problem code:} catalogue entry
\subsection{Definitions and mathematical statement}
The supplied record concerns a historical research question rather than a specified mathematical proposition. The expression ``mixed mathematics'' names the historical category under investigation; the record supplies neither a definition of that category nor criteria for assessing its relevance to mathematical practice. No formal definition can therefore be supplied without adding interpretive assumptions.

\emph{Source question.} Determine how relevant the category of mixed mathematics is to sixteenth-century mathematical practice, especially where the term is absent from the surviving sources.

The question includes practices documented in surviving sources that do not use the term ``mixed mathematics.'' The record does not specify a geographical scope, a corpus of surviving sources, the practices to be compared, or a criterion for relevance. Consequently, it does not determine a precise mathematical statement with explicit domains, hypotheses, and quantifiers. These unresolved historical and methodological qualifications require review before any more specific formulation is adopted.
\subsection{Short English statement}
Assess the historical relevance of ``mixed mathematics'' to sixteenth-century mathematical practice, giving particular attention to sources in which the term does not occur. The supplied question leaves the category, evidence base, and criteria for relevance unspecified, so a precise mathematical proposition cannot be extracted without additional assumptions.
\subsection{Sources}
\begin{itemize}
\item Local UnsolvedMath registry, permanent ID \(30002308\), title ``Mixed Mathematics in Sixteenth-Century Practice,'' source problem code \texttt{catalogue entry}, in the collection ``Oberwolfach Reports Open Problems.'' The supplied record gives 2013 as the proposed year and identifies no individual proposer.
\item Supplied source citation: \emph{From ``Mixed'' to ``Applied'' Mathematics: Tracing an important dimension of mathematics and its history} (2013), Oberwolfach Reports, pp.~657--733, DOI: \url{https://doi.org/10.4171/owr/2013/12}. The background describes Oberwolfach Reports as documenting mathematical workshops and their open-problem discussions. The supplied literature-source label is ``Oberwolfach Report 12/2013,'' with the same link: \url{https://doi.org/10.4171/owr/2013/12}.
\item Supplied dated literature curation, checked \texttt{2026-08-21}: the record labels the question ``open'' and reports that its resolution depends on unresolved historical comparisons of disciplinary practice and terminology, with no definitive synthesis found. This is supplied dated curation derived from the OWR record, not an independently verified literature assessment; time-sensitive status and cited claims require independent verification.
\item The supplied record marks the statement status as ``exact'' and reports that its mathematical content was checked against the cited official report, with notation, context, and question form normalized without changing the source-level claim. That verification is a claim of the supplied record and has not been independently checked here.
\end{itemize}
\end{document}
